This chapter reports the results of the study. Unless stated otherwise, every
number is measured on the held-out ISIC-2019 test split (4{,}796 images) at the
eight-view test-time-augmentation (TTA) setting, using the checkpoint selected by
the best validation BMA. All values are from a single run with seed 42.

\section{Performance Evaluation}

Table~\ref{tab:main} compares the two proposed models against the five
lightweight baselines, with every metric reported at the same TTA inference
setting so the rows are directly comparable. The lightweight hybrid reaches a BMA
of 0.6081 with only 3.98M parameters, which is the best balanced accuracy of any
model in the sub-6M-parameter budget ahead of MobileViT-S (0.6020), its
closest competitor, at fewer parameters and about two-thirds of the compute.
The flagship DEKAN model reaches 0.6438 BMA among the main proposed models and
also leads on every secondary metric (macro-F1 0.6059, macro-AUC 0.9271,
accuracy 0.7467) at the cost of roughly four times the parameters and ten times
the compute of the lightweight model. A linear-head ablation variant of DEKAN
(dekan\_linear) reaches 0.6510 BMA, which is the highest single-model figure in
the study; the contribution of the KAN classifier relative to a plain linear head
is analysed in Section~\ref{sec:ablations}. Figure~\ref{fig:comparison} shows
the same comparison graphically.

\begin{table}[H]
\centering
\caption{Main results on the ISIC-2019 test set (with 8-view TTA). Best value in
each column is in \textbf{bold}; the two proposed models are marked. A
linear-head ablation of DEKAN (dekan\_linear) reaches a higher BMA of 0.6510
but is reported with the other ablations in Section~\ref{subsec:ablation-study}
rather than here, so this table is not its intended place; the bold BMA marks
the best main model.}
\label{tab:main}
\resizebox{\textwidth}{!}{%
\begin{tabular}{lrrrrrr}
\toprule
Model & Params (M) & GMACs & BMA & Macro-F1 & Macro-AUC & Accuracy \\
\midrule
ResNet-18            & 11.69 & 1.810 & 0.3368 & 0.2763 & 0.8217 & 0.5225 \\
MobileNetV2          & 3.50  & 0.300 & 0.4738 & 0.4701 & 0.8827 & 0.6731 \\
EfficientNet-B0      & 5.29  & 0.390 & 0.4946 & 0.4736 & 0.9037 & 0.6779 \\
MobileViT-S          & 5.60  & 1.000 & 0.6020 & 0.5354 & 0.9218 & 0.6981 \\
EfficientFormer-L1   & 12.27 & 1.300 & 0.5519 & 0.5779 & 0.9005 & 0.7333 \\
\midrule
Hybrid (ours, light) & 3.98  & 0.631 & 0.6081 & 0.5602 & 0.9229 & 0.7162 \\
DEKAN (ours, flagship) & 16.45 & 6.633 & \textbf{0.6438} & \textbf{0.6059} & \textbf{0.9271} & \textbf{0.7467} \\
\bottomrule
\end{tabular}
}
\end{table}

\begin{figure}[H]
\centering
\includegraphics[width=\textwidth]{02_model_comparison.pdf}
\caption{Model comparison on the ISIC-2019 test set: BMA (with and without TTA)
on the left, macro-F1 and macro-AUC on the right. The two proposed models are
highlighted.}
\label{fig:comparison}
\end{figure}

The training behaviour of the lightweight hybrid is shown in
Figure~\ref{fig:curves}. Training and validation loss decrease smoothly and the
validation BMA peaks at 0.6505 at epoch 149, with no sign of the rare-class
collapse that a poorly tuned imbalance loss would cause evidence that the
$\beta=0.999$ choice in the Class-Balanced Focal Loss was appropriate.

\begin{figure}[H]
\centering
\includegraphics[width=\textwidth]{01_training_curves.pdf}
\caption{Training and validation loss (left) and BMA (right) for the lightweight
hybrid model over 150 epochs. The dashed line marks the best validation BMA.}
\label{fig:curves}
\end{figure}

DEKAN follows a similar pattern, shown in Figure~\ref{fig:curvesdekan}. Training
and validation loss decrease smoothly over the full 150 epochs, and the
validation BMA peaks late in training, again with no sign of rare-class
collapse consistent with the loss configuration transferring well from the
lightweight model to the higher-capacity dual-backbone setting.

\begin{figure}[H]
\centering
\includegraphics[width=\textwidth]{09_training_curves_dekan.pdf}
\caption{Training and validation loss (left) and BMA (right) for the DEKAN
model over 150 epochs. The dashed line marks the best validation BMA.}
\label{fig:curvesdekan}
\end{figure}

\section{Final Design Refinement}

The final architectures presented in Chapter 4 are the direct result of iterative refinement, driven by extensive evaluation and ablation analysis. Several structural and procedural modifications were introduced throughout development to systematically address technical vulnerabilities identified during preliminary baseline experiments.

Early experimental trials revealed that severe class imbalance caused highly unstable learning behavior, which manifested as poor recognition of minority classes. To prevent rare-class collapse, the training framework was upgraded to incorporate Class-Balanced Focal Loss paired with a weighted sampler. This modification significantly stabilized optimization and improved balanced multi-class accuracy. Similarly, handling missing metadata posed an early challenge; conventional preprocessing methods required statistical imputation, which inadvertently introduced fabricated clinical values into the training loop. To preserve data integrity, imputation was replaced with a learned ``missing'' embedding strategy, enabling the network to explicitly represent absent information without generating artificial data.

Training stability also emerged as a significant barrier when optimization involved transformer-based components. To curb overfitting and enhance generalization on the constrained dataset, we integrated stochastic depth, weight Exponential Moving Average (EMA), Mixup augmentation, and RandAugment into the final execution pipeline. 

System efficiency and OS-specific constraints required further engineering refinements. The Windows operating environment imposed strict synchronization limits on multi-worker data loading, creating a severe bottleneck that left the GPU underutilized. We bypassed this restriction entirely by building a custom memory-mapped dataset architecture coupled with a zero-worker loading strategy, ensuring high-throughput data delivery and strict reproducibility within our hardware constraints. 

Finally, ablation experiments indicated that the much-hyped KAN classifier head did not consistently outperform a standard linear layer. Rather than filtering out this negative outcome, both variants were retained and documented to maintain empirical transparency. 

Collectively, these refinements yield an architecture that balances diagnostic accuracy, computational efficiency, algorithmic fairness, and deployment readiness.

\begin{table}[H]
\centering
\caption{Summary of empirical issues identified during development and the corresponding design modifications implemented in the final architecture.}
\label{tab:design_refinements}
\begin{tabularx}{\textwidth}{>{\raggedright\arraybackslash}p{0.38\textwidth}>{\raggedright\arraybackslash}X}
\toprule
\textbf{Identified Issue} & \textbf{Design Modification} \\
\midrule
Rare-class collapse under imbalance & Class-Balanced Focal Loss and weighted sampling \\
Missing metadata values & Learned ``missing'' embeddings \\
Transformer instability and overfitting & Stochastic depth, Mixup, RandAugment, and weight EMA \\
Slow data loading / Windows deadlock & Memory-mapped dataset and zero-worker loading \\
Uncertain KAN head effectiveness & Linear-head ablation retained and reported transparently \\
\bottomrule
\end{tabularx}
\end{table}


\section{Per-Class Analysis}

Because BMA is an average over classes, it is worth looking at where the models
succeed and fail per class. Table~\ref{tab:perclass} lists per-class recall for
two representative baselines and the two proposed models, and
Figure~\ref{fig:perclass} compares the lightweight hybrid against EfficientNet-B0
class by class.

\begin{table}[H]
\centering
\caption{Per-class recall on the ISIC-2019 test set (with TTA).}
\label{tab:perclass}
\begin{tabular}{lrrrr}
\toprule
Class & EfficientNet-B0 & MobileViT-S & Hybrid (ours) & DEKAN (ours) \\
\midrule
MEL  & 0.594 & 0.447 & 0.518 & 0.611 \\
NV   & 0.857 & 0.890 & 0.859 & 0.859 \\
BCC  & 0.546 & 0.467 & 0.618 & 0.731 \\
AK   & 0.348 & 0.452 & 0.587 & 0.368 \\
BKL  & 0.366 & 0.587 & 0.632 & 0.622 \\
DF   & 0.245 & 0.660 & 0.566 & 0.509 \\
VASC & 0.843 & 0.882 & 0.863 & 0.922 \\
SCC  & 0.157 & 0.430 & 0.223 & 0.529 \\
\bottomrule
\end{tabular}
\end{table}

The pattern is consistent with the headline numbers. The lightweight hybrid lifts
several of the hard middle classes relative to the convolutional baselines
BCC, AK, and BKL in particular which is where its global reasoning and
metadata context help most. Notably, the hybrid's AK recall (0.587) is the
highest of any model in the study, including DEKAN (0.368); actinic keratosis
is a sun-damage-related precancerous condition disproportionately common in
older patients at exposed anatomical sites, precisely the profile where the
metadata cross-attention shifts the prior most reliably. Its weakest class is
SCC (0.223), a rare class that is easily confused with AK and BKL. DEKAN spreads its recall more evenly and
recovers SCC substantially (0.529) while also posting the strongest VASC (0.922)
and BCC (0.731), which is the main reason its balanced accuracy is highest.
Figure~\ref{fig:confusion} shows the full confusion matrix for the lightweight
hybrid; the off-diagonal mass sits mostly among the keratinocytic classes
(AK/BKL/SCC), exactly the clinically difficult group.

\begin{figure}[H]
\centering
\includegraphics[width=0.85\textwidth]{04_per_class_recall.pdf}
\caption{Per-class recall of the lightweight hybrid versus EfficientNet-B0,
sorted by the hybrid's recall.}
\label{fig:perclass}
\end{figure}

\begin{figure}[H]
\centering
\includegraphics[width=0.8\textwidth]{03_confusion_matrix.pdf}
\caption{Confusion matrix for the lightweight hybrid (normalised by true class
with TTA).}
\label{fig:confusion}
\end{figure}

Figure~\ref{fig:dekanconfusion} shows the same matrix for the DEKAN flagship.
The keratinocytic confusion (AK/BKL/SCC) is reduced relative to the lightweight
hybrid, and DEKAN recovers SCC substantially (0.529 vs.\ 0.223) the class
that costs the hybrid most. VASC recall reaches 0.922, the highest of any model
in the study.

\begin{figure}[H]
\centering
\includegraphics[width=0.8\textwidth]{08_confusion_matrix_dekan.pdf}
\caption{Confusion matrix for the DEKAN flagship model (normalised by true class,
with TTA).}
\label{fig:dekanconfusion}
\end{figure}

\section{Efficiency Analysis}

The point of the lightweight model is accuracy per unit of cost so
Table~\ref{tab:efficiency} reports balanced accuracy per million parameters
alongside the raw numbers and Figure~\ref{fig:efficiency} plots accuracy against
parameter count with compute shown as bubble size. Among the full baselines and proposed
models, the lightweight hybrid has the best BMA-per-parameter (0.153). The
isolated CNN-stem (0.268) and pure-ViT (0.249) ablations score higher on this
ratio alone, but both sit far below the full hybrid in absolute BMA, so their
apparent efficiency is misleading parameter-efficiency without sufficient
accuracy is not a useful operating point. The hybrid is therefore the best
operating point that also clears a useful accuracy level. DEKAN sits at the opposite end: it buys the
highest absolute accuracy with the lowest efficiency, which is the expected and
intended trade-off between the two tiers.

\begin{table}[H]
\centering
\caption{Efficiency: balanced accuracy per million parameters (TTA BMA).}
\label{tab:efficiency}
\begin{tabular}{lrrrr}
\toprule
Model & Params (M) & GMACs & BMA & BMA / M-param \\
\midrule
ResNet-18            & 11.69 & 1.810 & 0.3368 & 0.029 \\
MobileNetV2          & 3.50  & 0.300 & 0.4738 & 0.135 \\
EfficientNet-B0      & 5.29  & 0.390 & 0.4946 & 0.093 \\
MobileViT-S          & 5.60  & 1.000 & 0.6020 & 0.108 \\
EfficientFormer-L1   & 12.27 & 1.300 & 0.5519 & 0.045 \\
Hybrid (ours, light) & 3.98  & 0.631 & 0.6081 & \textbf{0.153} \\
DEKAN (ours, flagship) & 16.45 & 6.633 & 0.6438 & 0.039 \\
\bottomrule
\end{tabular}
\end{table}

\begin{figure}[H]
\centering
\includegraphics[width=0.85\textwidth]{05_efficiency_scatter.pdf}
\caption{Accuracy vs.\ parameters; bubble size is proportional to GMACs. The
dotted line marks the 6M-parameter budget. The lightweight hybrid (red) gives the
best accuracy within budget.}
\label{fig:efficiency}
\end{figure}

\section{Effect of Test-Time Augmentation}

TTA helps most but not all models, so Table~\ref{tab:tta} reports BMA with and
without it. The lightweight hybrid (full model) benefits most among the strong
models ($+0.027$), and DEKAN gains a further $+0.007$. Among baselines,
EfficientFormer-L1 drops slightly and ResNet-18 degrades sharply ($-0.151$);
the ResNet-18 case is a known failure mode in which per-view predictions are
confidently wrong in different directions, so averaging hurts.

The no-metadata ablation variant reveals a further pattern. Without the metadata
head, the hybrid goes from 0.6216 to 0.6085 under TTA ($-0.013$) the same
degradation direction as ResNet-18. Adding the metadata head reverses this:
the full model goes from 0.5815 to 0.6081 ($+0.027$). The two variants
converge to nearly the same TTA number (0.6085 vs 0.6081) from opposite sides.
This TTA$\times$metadata interaction suggests the cross-attention head
stabilises per-view predictions; without it the model becomes overconfident
in different directions across augmented views.

\begin{table}[H]
\centering
\caption{BMA without and with 8-view TTA, and the change.}
\label{tab:tta}
\begin{tabular}{lrrr}
\toprule
Model & BMA (no TTA) & BMA (TTA) & $\Delta$ \\
\midrule
ResNet-18            & 0.4875 & 0.3368 & $-0.151$ \\
MobileNetV2          & 0.4485 & 0.4738 & $+0.025$ \\
EfficientNet-B0      & 0.4917 & 0.4946 & $+0.003$ \\
MobileViT-S          & 0.5925 & 0.6020 & $+0.010$ \\
EfficientFormer-L1   & 0.5559 & 0.5519 & $-0.004$ \\
\midrule
Hybrid w/o metadata  & 0.6216 & 0.6085 & $-0.013$ \\
Hybrid (ours, light) & 0.5815 & 0.6081 & $+0.027$ \\
DEKAN (ours, flagship) & 0.6368 & 0.6438 & $+0.007$ \\
\bottomrule
\end{tabular}
\end{table}

\section{Analysis of Design Solutions (Ablations)}
\label{subsec:ablation-study}

\subsection{Lightweight hybrid}

Table~\ref{tab:ablhybrid} and Figure~\ref{fig:ablation} isolate the components of
the lightweight model. The pure-ViT ablation (no CNN stem, 1.97M parameters)
reaches 0.4908 BMA, virtually identical to the CNN-stem-only result (0.4853). The
two components are roughly equal in isolation but strongly complementary when
combined: adding the ViT trunk to the CNN stem yields a gain of $+0.117$ to
$+0.123$ BMA, and the full hybrid reaches 0.6081. This is the primary
architectural justification for the hybrid design neither local nor global
features alone approach the performance the model achieves when both are present.

The effect of metadata in isolation is however modest at test time on this
single run: the full model (0.6081) and the no-metadata variant (0.6085) are
within 0.0004 of each other on the test set even though on the validation set
the full model led by about 0.011 BMA. We therefore report the metadata
contribution for the lightweight model honestly as inconclusive on a single test
evaluation; the DEKAN results below give a clearer signal. The no-TTA numbers add
a structural observation: without metadata the model scores 0.6216 before
augmentation (versus 0.5815 for the full model), a gap of 0.040 that is entirely
closed by TTA the full model gains $+0.027$ while the no-metadata variant
loses $-0.013$ (Table~\ref{tab:tta}). The two variants converge to the same TTA
result from opposite sides, which points to the metadata head regularising
per-view predictions rather than lifting the mean class probability.

\begin{table}[H]
\centering
\caption{Lightweight hybrid ablation (TTA BMA). All four variants were trained
and evaluated.}
\label{tab:ablhybrid}
\resizebox{\textwidth}{!}{%
\begin{tabular}{llr}
\toprule
Configuration & Variant & BMA (TTA) \\
\midrule
CNN stem + GAP + linear     & hybrid\_cnn\_only & 0.4853 \\
Pure ViT, no CNN stem       & hybrid\_vit\_only & 0.4908 \\
CNN + ViT, no metadata      & hybrid\_no\_meta  & 0.6085 \\
CNN + ViT + metadata (full) & hybrid\_full      & 0.6081 \\
\bottomrule
\end{tabular}
}
\end{table}

\begin{figure}[H]
\centering
\includegraphics[width=0.85\textwidth]{06_ablation.pdf}
\caption{Lightweight hybrid ablation: BMA for all four variants. CNN-only and
ViT-only perform similarly in isolation; the full hybrid is substantially higher.}
\label{fig:ablation}
\end{figure}

\subsection{DEKAN}
\label{sec:ablations}

Table~\ref{tab:abldekan} and Figure~\ref{fig:dekanablation} isolate two effects
in the flagship model: the contribution of clinical metadata and the choice of
classifier head (KAN vs.\ linear).

\textbf{Metadata.} Comparing dekan\_no\_meta (0.6302) with dekan\_linear
(0.6510) both use a linear classifier head, differing only in the presence of
the metadata cross-attention block shows a gain of $+0.021$ BMA. This is a
clear and consistent signal that patient age, sex, and anatomical site provide
useful diagnostic context at this model capacity.

\textbf{KAN vs.\ linear head.} Replacing the KAN classifier with a plain linear
layer (dekan\_linear, 0.6510) while keeping everything else identical produces a
higher BMA than the full KAN model (dekan\_full, 0.6438), a difference of
$+0.007$ in favour of the linear head. The KAN classifier therefore does not
improve over a plain linear layer on this task. This is an honest negative result:
the KAN's learnable splines add no measurable value at the classifier stage of a
well-trained vision backbone on this noisy imbalanced 8-class dataset. The
likely explanation is that the discriminative difficulty on ISIC-2019 lies in the
feature representation, not in the shape of the final decision boundary which a
linear layer can already model adequately after the rich dual-backbone
representation.

\begin{table}[H]
\centering
\caption{DEKAN ablation: effect of clinical metadata and the KAN vs.\ linear
classifier head (TTA BMA). The single-backbone variants were not trained, so
the dual-backbone fusion gain cannot yet be quantified.}
\label{tab:abldekan}
\resizebox{\textwidth}{!}{%
\begin{tabular}{llr}
\toprule
Configuration & Variant & BMA (TTA) \\
\midrule
Dual backbone + fusion, no metadata, linear head & dekan\_no\_meta  & 0.6302 \\
Dual backbone + fusion + metadata, linear head   & dekan\_linear    & 0.6510 \\
Dual backbone + fusion + metadata + KAN (full)   & dekan\_full      & 0.6438 \\
\bottomrule
\end{tabular}
}
\end{table}

\begin{figure}[H]
\centering
\includegraphics[width=\textwidth]{07_dekan_ablation.pdf}
\caption{DEKAN ablation: effect of clinical metadata and KAN vs.\ linear
classifier head on BMA. Single-backbone variants were not trained and are absent.}
\label{fig:dekanablation}
\end{figure}

\section{Comparisons and Relationships}

Read against the wider literature, these results sit where one would expect a
single-model, efficiency-constrained system to sit. The top entries on the
ISIC-2019 challenge leaderboard~\cite{isic2019leaderboard} reach a balanced
accuracy around 0.636 (DAISYLab) using large ensembles of multi-resolution
EfficientNets and SENet-154. DEKAN, a single model, reaches 0.6438 BMA on our
held-out split, and even the 3.98M-parameter lightweight hybrid reaches 0.6081
within striking distance of an entire ensemble at a tiny fraction of the cost.
The comparison is not exact since the leaderboard is computed on the official
challenge test server and our split is an internal one, so we present it as
context rather than as a like-for-like ranking. Relative to controlled
single-model comparisons on this dataset such as Aktan et
al.~\cite{aktan2025comparison}, our contribution is to add clinical-metadata
fusion to the comparison and to report accuracy together with a hard efficiency
budget. The clearest relationship to prior work is with the metadata-fusion line
of Pacheco and Krohling~\cite{pacheco2020impact,pacheco2021attention}: we adopt
their core idea of representing missing values with learned embeddings, and we
extend their heavy-CNN setting to lightweight and dual-backbone hybrids.

\section{Discussions}

The results support the central claims of the thesis while leaving a few
questions open. The lightweight hybrid does what it was designed to do: it beats
all five same-budget baselines on balanced accuracy and is the most
parameter-efficient model in its class which makes it a credible candidate for
deployment on modest hardware. DEKAN shows that the same design scales: investing
capacity in a second backbone, learned fusion and a deeper trunk pushes balanced
accuracy to the level of a leaderboard ensemble in a single model. Metadata
fusion helps clearly in the higher-capacity model and is at worst harmless in
the lightweight one.

Two findings cut against initial expectations and are worth stating plainly.
First, the benefit of metadata for the lightweight model could not be confirmed
on a single test evaluation, even though the validation curve favoured it by
0.011 BMA. The DEKAN results give a clearer signal at higher capacity
($+0.021$ BMA between dekan\_no\_meta and dekan\_linear), suggesting the
metadata effect requires sufficient model capacity to manifest consistently.
Second, the KAN classifier in DEKAN does not outperform a plain linear head:
dekan\_linear (0.6510) beats dekan\_full (0.6438) by 0.007. This is the
highest single-model BMA achieved in the study, but it belongs to the
linear-head ablation, not to the KAN model, which is the intended flagship.
The KAN's expressive spline activations appear to offer no advantage over a
linear classifier once the dual-backbone representation is rich enough.

\section{Limitations}

Several limitations bound these conclusions. All numbers come from a single seed
(42); publication-grade claims would require repeating each run over multiple
seeds and reporting mean and standard deviation which we have not yet done.
Two ablation variants were not trained dekan\_densenet\_only and
dekan\_effnet\_only so the individual contribution of each backbone and the
precise size of the dual-backbone fusion gain remain unquantified. We also did
not evaluate on an external dataset so the generalisation of these models beyond
the ISIC-2019 distribution is untested. Finally, the comparison with the
ISIC-2019 leaderboard is contextual not a controlled head-to-head because the
evaluation splits differ.
